\documentclass[../../main.tex]{subfiles}
\begin{document}

\begin{flushright}
\footnotesize\textit{【自動文字起こし・要確認】}
\end{flushright}

[解] 題意から$y = \dfrac{a}{x} + b$の\\
グラフは右のようである\\
（対称性から$a>0$のみ考える）\\
よって、$y = \dfrac{a}{x} + b$ が$(1, 8), (5, 0)$を通る\\
から
\[
\begin{cases}
8 = a + b \\
0 = \dfrac{a}{5} + b
\end{cases}
\Leftrightarrow
\begin{cases}
b = -2 \\
a = 10
\end{cases}
\]
である。以下 $f(x) = \dfrac{10}{x} - 2$ とおく。

\begin{center}
\begin{tikzpicture}[scale=0.5]
    \draw[->] (-1, 0) -- (7, 0) node[right] {$x$};
    \draw[->] (0, -3) -- (0, 9) node[above] {$y$};
    \draw[domain=1:6, smooth, variable=\x] plot ({\x}, {10/\x - 2});
    \draw[dashed] (1, 0) -- (1, 8) -- (0, 8) node[left] {$8$};
    \node[below] at (5, 0) {$5$};
\end{tikzpicture}
\end{center}

\paragraph{(1)}

高さ$h$ ($0 \le h \le 8$) の時の水の体積\\
$V(h)$とおくと
\begin{align}
V(h) &= \int_0^h \pi \left( \dfrac{10}{y+2} \right)^2 dy \\
&= 100\pi \left[ -\dfrac{1}{y+2} \right]_0^h \\
&= 100\pi \left( \dfrac{1}{2} - \dfrac{1}{h+2} \right)
\end{align}
だから
\[ V(6) = \dfrac{75}{2}\pi \quad // \]

\paragraph{(2)}

$\left. \dfrac{dh}{dt} \right|_{h=3}$ をもとめれば良い。題意から $\dfrac{dV}{dt} = k$ であって、
$\dfrac{dh}{dt} = h'$ とおくと、
\[ k = 100\pi \dfrac{1}{(h+2)^2} h' \]
だから$h=3$として、
\[ \left. h' \right|_{h=3} = \dfrac{25}{100\pi}k = \dfrac{k}{4\pi} \text{ [cm/s]} \quad // \]

\vspace{1cm}
[解 2(2)] エースを使ってみる。\\
\paragraph{(2)}

（$y = \dfrac{10}{x} - 2$まで同じ）時刻$t$での表面積$S$、高さ$h$とすると
\[ k = S \dfrac{dh}{dt} \quad \cdots \text{①} \]
\[ S = \pi \left(\dfrac{10}{h+2}\right)^2 \quad \cdots \text{②} \]
②を①に代入
\[ \dfrac{dh}{dt} = \dfrac{k}{\left(\dfrac{10}{h+2}\right)^2 \cdot \pi} \]
$h=3$として
\[ \dfrac{dh}{dt} = \dfrac{k}{4\pi} \quad // \]

{\color{cyan}
\fbox{エースのすばらしさ!! これはとんでもない}
}

\end{document}
